\section{Separated homogeneous peaks}\label{sec:peaks}

We begin the construction of the potential in
Proposition~\ref{base:existence}.  Its analytic requirement will follow
from increasingly negative averages against positive measures that
preserve holomorphic means.  The following observation explains the
role of those measures before we construct them.

Write $S=\{\xi\in\mathbb C^2:|\xi|=1\}$, and let $\sigma$ be its
surface measure normalized to have mass one.  A probability measure
$\mu$ on $S$ \emph{represents holomorphic evaluation at zero} if
\[
 \int_S p(\xi)\,d\mu(\xi)=p(0)
 \quad\text{for every holomorphic polynomial }p.
\]

\begin{lemma}[Representing-mean criterion]\label{peaks:mean-criterion}
Let $u$ be a continuous real function on $\mathbb B$.  Suppose there
are radii $0<r_N<1$ and representing probabilities $\mu_N$ on $S$
such that
\[
 \int_S u(r_N\xi)\,d\mu_N(\xi)\longrightarrow-\infty.
\]
Then no holomorphic function $H$ on $\mathbb B$ and constant $C_H$
satisfy $\operatorname{Re}H\le C_H+u$ on $\mathbb B$.
\end{lemma}

\begin{proof}
For each fixed $r_N<1$, the Taylor polynomials of $H$ converge
uniformly on $r_NS$.  The representing identities therefore give
\[
 \int_S\operatorname{Re}H(r_N\xi)\,d\mu_N(\xi)
 =\operatorname{Re}H(0).
\]
Integrating the proposed inequality contradicts the assumed limit.
\end{proof}

We will apply this criterion to
$u(z)=\phi(z)+4\log(1/(1-|z|))$.  The first ingredient is a supply
of homogeneous polynomials with many separated peaks.  Their uniform
bounds will also control the derivatives of the eventual potential.

\subsection{Projective packing and polynomial bounds}

The projective distance and separated-power estimate used here occur
in the direct construction of Ryll and Wojtaszczyk
\cite[p.~112 and Lemma~2.7]{RyllWojtaszczyk1983}.
We give the packing argument with a Gaussian tail and an adjustable
spacing parameter, since the density corrections will require both.
Related constructions with peaks localized in nonisotropic sphere
balls appear in Aleksandrov \cite[Lemmas~2--3]{Aleksandrov1983};
the caps used here are invariant under a common phase rotation.

Identify complex lines in $\mathbb C^2$ with points of
$\mathbb {CP}^1$.  For unit representatives define
\[
 d([\xi],[p])=\sqrt{1-|\langle\xi,p\rangle|^2},
\]
where the Hermitian inner product is linear in its first argument.
This equals the operator norm distance between the rank-one
orthogonal projections onto the two lines: their difference has
trace zero and eigenvalues
$\pm\sqrt{1-|\langle\xi,p\rangle|^2}$.  The triangle inequality
for the operator norm therefore proves that $d$ is a metric.  We also write
$d(\xi,p)$ when the representatives are understood.

In unitary coordinates with $\xi_1=\langle\xi,p\rangle$, write
\[
 \xi=(e^{i\theta}\cos v,e^{i\gamma}\sin v),
 \qquad 0\le v\le\pi/2.
\]
The round volume element is
$\cos v\sin v\,dv\,d\theta\,d\gamma$.  After normalization and
the change of variable $u=\cos^2v$, this gives
\begin{equation}\label{peaks:sphere-measure}
 d\sigma=du\,\frac{d\theta}{2\pi}\frac{d\gamma}{2\pi},
 \qquad 0\le u\le1.
\end{equation}
In particular, a projective cap $d(\xi,p)<s$ has measure $s^2$
for $0\le s\le1$.

Fix $D_0\ge1$.  For an integer $k\ge D_0^2$, choose a maximal
$D_0/\sqrt{k}$-separated set of lines, and choose one unit
representative $p$ for each line.  Let $\mathcal P_k$ denote this
set of representatives and $N_k=|\mathcal P_k|$.  Define the
homogeneous holomorphic polynomial
\begin{equation}\label{peaks:polynomial}
 P_k(z)=\sum_{p\in\mathcal P_k}\langle z,p\rangle^k.
\end{equation}
The representatives may be chosen arbitrarily; the estimates below
hold for every such choice.

\begin{lemma}[Peak bounds]\label{peaks:bounds}
There are absolute constants $C,c>0$, independent of $D_0\ge1$,
such that
\begin{equation}\label{peaks:count}
 \frac{k}{D_0^2}\le N_k\le\frac{4k}{D_0^2}
\end{equation}
and, for every $\xi\in S$ and $u\ge0$,
\begin{equation}\label{peaks:gaussian}
 \sum_{\substack{p\in\mathcal P_k\\\sqrt{k}d(\xi,p)\ge u}}
       |\langle\xi,p\rangle|^k\le Ce^{-cu^2}.
\end{equation}
Consequently,
\begin{equation}\label{peaks:global}
 |P_k(z)|\le C|z|^k\quad(z\in\mathbb C^2),
 \qquad
 \int_S|P_k|\,d\sigma\le\frac{2N_k}{k+2}.
\end{equation}
\end{lemma}

\begin{proof}
Put $\rho=D_0/\sqrt{k}\le1$.  Maximality gives a covering by caps
of radius $\rho$, and separation gives disjoint interiors for caps
of radius $\rho/2$.  Their measures yield
$1\le N_k\rho^2$ and $N_k\rho^2/4\le1$, proving
\eqref{peaks:count}.

We also need a local count whose constant does not depend on $D_0$.
About the centers within distance $u/\sqrt{k}$ of a fixed line,
place caps of radius $1/(3\sqrt{k})$.  Their interiors are disjoint,
and they lie inside the cap of radius $(u+1/3)/\sqrt{k}$ about that
line.  Comparing measures bounds their number by $9(u+1/3)^2$ when
the latter radius is at most one.  Otherwise their number is at
most $9k\le9(u+1/3)^2$, by comparison with the whole space.
Thus the local count is at most $C(1+u)^2$ in all cases.

Since
\[
 |\langle\xi,p\rangle|^k
 =(1-d(\xi,p)^2)^{k/2}
 \le e^{-kd(\xi,p)^2/2},
\]
the contribution from the annulus
$n\le\sqrt{k}d(\xi,p)<n+1$ is at most
$C(n+2)^2e^{-n^2/2}$.  Summing these bounds from $n=\lfloor u\rfloor$
onward proves \eqref{peaks:gaussian}: the polynomial factor is
absorbed by a weaker Gaussian exponent, and the range $0\le u\le2$
is absorbed by the constant.  Setting $u=0$ and using homogeneity
gives the first assertion in \eqref{peaks:global}, including points
outside the unit ball.  Finally, \eqref{peaks:sphere-measure} gives
\[
 \int_S|\langle\xi,p\rangle|^k\,d\sigma
 =\int_0^1 u^{k/2}\,du=\frac2{k+2}.
\]
The triangle inequality proves the stated integral bound.
\end{proof}

\subsection{Patches and a radial profile}

For a fixed $R_0>0$ and a center $p\in\mathcal P_k$, use the
adapted coordinates above and put
\begin{equation}\label{peaks:patch-definition}
 y=-k\log|\xi_1|,\qquad
 \mathcal U_{k,p}=\{\xi\in S:0\le y<R_0\}.
\end{equation}
The set $y=0$ is the phase circle through $p$.  Every patch is
invariant under multiplication of $\xi$ by a common phase.

\begin{lemma}[Patch estimates]\label{peaks:patches}
Fix $R_0>0$.  If $D_0>\max\{1,2\sqrt{2R_0}\}$ and
$k\ge\max\{D_0^2,2R_0\}$, the patches $\mathcal U_{k,p}$
are pairwise disjoint.  On each patch,
\begin{equation}\label{peaks:patch-measure}
 d\sigma=\frac2k e^{-2y/k}\,dy\,
             \frac{d\theta}{2\pi}\frac{d\gamma}{2\pi},
\end{equation}
and
\begin{equation}\label{peaks:single-peak}
 |P_k(\xi)-e^{-y}e^{ik\theta}|
 \le C e^{-c(D_0-\sqrt{2R_0})^2}.
\end{equation}
Outside their union,
\begin{equation}\label{peaks:off-patch}
 |P_k(\xi)|\le Ce^{-cR_0}.
\end{equation}
The constants $C,c$ can be chosen independently of $R_0,D_0,k$.
\end{lemma}

\begin{proof}
The projective radius of a patch is
$\sqrt{1-e^{-2R_0/k}}\le\sqrt{2R_0/k}$, so the separation of
the centers makes the patches disjoint.  Formula
\eqref{peaks:patch-measure} follows from
$|\xi_1|^2=e^{-2y/k}$ in \eqref{peaks:sphere-measure}.

If $\xi\in\mathcal U_{k,p}$ and $p'\ne p$ is another center,
the metric triangle inequality gives
\[
 \sqrt{k}d(\xi,p')\ge D_0-\sqrt{2R_0}.
\]
The peak from $p$ is exactly $e^{-y}e^{ik\theta}$.
Summing all other peaks by \eqref{peaks:gaussian} proves
\eqref{peaks:single-peak}.
Outside every patch, each center satisfies
\[
 kd(\xi,p)^2\ge k(1-e^{-2R_0/k})\ge R_0.
\]
The final inequality uses $k\ge2R_0$ and
$1-e^{-v}\ge v/2$ for $0\le v\le1$.  Another application of
\eqref{peaks:gaussian} proves \eqref{peaks:off-patch}.
\end{proof}

The density correction will use a radial profile with an exact
weighted cancellation.  The following choice also supplies the
strict inequality needed for the logarithmic tests later.

\begin{lemma}[Radial profiles]\label{peaks:profiles}
Fix $1<a<2$ and put $\alpha=\log a$.  One can choose $0<s_1<1$,
smooth real functions $f,b$ on $[0,\infty)$, and $R_0>0$ such that
\begin{equation}\label{peaks:profile-negativity}
 \alpha^2-2s_1(\cosh\alpha-1)<0,
\end{equation}
\begin{equation}\label{peaks:profile-identities}
 \|f\|_\infty<1,\quad f=-s_1\ \text{on }[0,\alpha],\quad
 \int_0^\infty e^{-y}f(y)\,dy=0,\quad
 b'-b=f,\quad b(0)=0.
\end{equation}
Both functions vanish on a neighborhood of $[R_0,\infty)$.
Moreover, $R_0$ can be chosen so that $|aP_k|\le1/2$ off the
patches whenever the hypotheses of Lemma~\ref{peaks:patches} hold.
\end{lemma}

\begin{proof}
The strict inequality $2(\cosh\alpha-1)>\alpha^2$ allows a choice
\[
 \frac{\alpha^2}{2(\cosh\alpha-1)}<s_1<1.
\]
Since $a<2$, this choice also satisfies
$s_1(1-e^{-\alpha})<e^{-\alpha}$.
Choose a small $\delta>0$ and a smooth cutoff $\chi$, equal to one
on $[0,\alpha]$, between zero and one, and zero for
$y\ge\alpha+\delta$.  For sufficiently small $\delta$,
\[
 N=s_1\int_0^\infty e^{-y}\chi(y)\,dy
 <e^{-(\alpha+2\delta)}.
\]
Choose a smooth compactly supported function $\nu$, supported
strictly after $\alpha+2\delta$, with $0\le\nu\le1$, and with
\[
 P=\int_0^\infty e^{-y}\nu(y)\,dy>N.
\]
Such a function is obtained by approximating the indicator of a
sufficiently long interval beginning just after $\alpha+2\delta$.
Then $f=-s_1\chi+(N/P)\nu$ has the required plateau and weighted
integral.  The two supports are disjoint, so
$\|f\|_\infty\le\max\{s_1,N/P\}<1$.

Define
\[
 b(y)=e^y\int_0^y e^{-v}f(v)\,dv.
\]
It satisfies $b'-b=f$ and $b(0)=0$.  Beyond the support of $f$,
the integral is exactly zero by the weighted cancellation, so $b$
also has compact support.  Near zero,
$b(y)=-s_1(e^y-1)$, which is smooth up to the boundary of the
half-line.  Smoothness here does not require extending $f$ by zero
to negative arguments.  Finally choose $R_0$ larger than both
supports, and enlarge it further until
$aCe^{-cR_0}\le1/2$ in \eqref{peaks:off-patch}.
\end{proof}

The plateau $0\le y\le\alpha$ is exactly the region where the
model peak $ae^{-y}$ has modulus at least one.  For the profiles
just constructed, the integral of
$\log|1-ae^{-y}e^{it'}|^2$ against
$(1+f(y)\cos t')\,dy\,dt'/(2\pi)$ on
$[0,R_0]\times[0,2\pi]$ equals
$\alpha^2-2s_1(\cosh\alpha-1)$, as calculated in
\eqref{tests:model-sign}.  Thus the strict profile inequality
prescribes a negative logarithmic integral, while the weighted
cancellation will preserve holomorphic moments.

Fix these profiles and $R_0$ from now on.  The spacing $D_0$ and
the degrees will remain available to be chosen sufficiently large.
In particular, enlarging $D_0$ later improves the single-peak
approximation without altering the profiles.
